\documentclass{style} % Use the custom style

\pretolerance=10000
\tolerance=2000

\hypersetup{
  colorlinks=true,
  urlcolor=black,
  linkcolor=black
}

\usepackage[left=0.4 in,top=0.4in,right=0.4 in,bottom=0.4in]{geometry} % Document margins
\newcommand{\tab}[1]{\hspace{.2667\textwidth}\rlap{#1}}
\newcommand{\itab}[1]{\hspace{0em}\rlap{#1}}
\name{Максимов Илья}

\address{
    \href{https://t.me/ilyhamaks}{t.me/ilyhamaks} \\
    \href{https://github.com/MaximovIlya}{github.com/MaximovIlya} \\
    \href{mailto:dilershina@gmail.com}{dilershina@gmail.com} 
}

\begin{document}

%----------------------------------------------------------------------------------------
%	OBJECTIVE
%----------------------------------------------------------------------------------------

%----------------------------------------------------------------------------------------
%	EDUCATION SECTION
%----------------------------------------------------------------------------------------

\begin{rSection}{Образование}

    {\bf Бакалавр}, Университет Иннополис \hfill {2023 - 2027}\\
    Искусственный Интеллект и Анализ Данных\\\\
    Факультет компьютерных и инженерных наук

\end{rSection}

%----------------------------------------------------------------------------------------
%	SKILLS SECTION
%----------------------------------------------------------------------------------------

\begin{rSection}{Навыки и умения}

\begin{itemize}
    \itemsep -3pt {}
     \item Языки программирования: Go, Python, Java, Dart, JavaScript
     \item Фреймворки: Flutter, Next.js, Node.js
     \item Инструменты: Docker, Git, CI/CD,  Swagger
     \item Базы Данных: PostgreSQL, SQLite
     \item Технологии: REST API, WebSocket, JWT, Clean Architecture, BloC
     \item Иностранный язык: английский (B2)
\end{itemize}

\end{rSection}

%----------------------------------------------------------------------------------------
%	PROJECTS SECTION
%----------------------------------------------------------------------------------------

\begin{rSection}{Проекты}

\begin{itemize}
    \itemsep -3pt {}
     \item \href{https://github.com/MaximovIlya/chat_app.git}{\textbf{Chat App}}: кросс-платформенный мессенджер (Android, iOS, Web, Desktop). Backend на Go/Fiber — REST API + WebSocket для real-time сообщений, JWT + bcrypt, PostgreSQL 17/pgx/sqlc, Swagger. Frontend на Flutter: Clean Architecture, BLoC, socket\_io\_client, Flutter SecureStorage.
     \item \href{https://github.com/MaximovIlya/tender_bot_go}{\textbf{Tender Bot}}: Telegram-бот на Go для управления тендерами по модели голландского аукциона. Три роли: организатор, поставщик, администратор. Автоматические уведомления через cron, деплой в Docker на Amvera.tech. Stack: telebot.v3, PostgreSQL/pgx/sqlc, robfig/cron.
     \item \href{https://github.com/OpenEventsMap/OpenEventsBackend}{\textbf{Open Events}}: платформа для поиска и создания публичных мероприятий на интерактивной карте. Backend на Go/Fiber — REST API + WebSocket-чат, JWT, CRUD событий с геокоординатами, управление участниками, cron-очистка завершённых событий. Stack: PostgreSQL 15/pgx/sqlc, golang-migrate, Swagger, Docker, Amvera.
     \item \href{https://github.com/IU-Capstone-Project-2025/Looki}{\textbf{Looki}}: мобильное приложение для виртуальной примерки одежды. Frontend на Flutter: авторизация, параметры тела, каталог коллекций, примерка одежды на 3D-модели, загрузка фото. 
\end{itemize}

\end{rSection}

%----------------------------------------------------------------------------------------
%	HACKATHONS SECTION
%----------------------------------------------------------------------------------------

\begin{rSection}{Практика}

   
    VK Practice \hfill \hfill Май 2026 - Июль 2026\
    \begin{itemize}
        \itemsep -3pt {}
        \item Разработал веб-приложение для проведения квизов в реальном времени: авторизация с ролями (организатор/участник), конструктор квизов, live-сессии с подключением по коду комнаты и лидербордом.
        \item Frontend: Next.js 14 (App Router), TypeScript, Tailwind CSS, Socket.IO client, NextAuth.js.
        \item Backend: Next.js API Routes, Socket.IO server, PostgreSQL, Prisma ORM.
        \item \href{https://github.com/MaximovIlya/vk_practice_project}{Репозиторий} \quad \href{https://drive.google.com/file/d/1UOm4LQo9R1axcW96MlJNprpl17DQTAPg/view?usp=sharing}{Сертификат}
    \end{itemize}

    Курс <<Системное программирование>>, VK Education \hfill 2026\\
    \begin{itemize}
        \itemsep -3pt {}
        \item Темы: ядро Linux, шедулинг процессов, процесс и его ресурсы, системные вызовы, устройство памяти, сигналы и прерывания, файловая система, потоки, IPC, сеть, пользователи, сессии, права доступа.
        \item \href{СЕРТИФИКАТ_ССЫЛКА}{Сертификат}
    \end{itemize}

\end{rSection}

\begin{rSection}{Хакатоны}
\vspace{-1.25em}
\item \textbf{Хакатон <<блокчейн разработка>> БАШНЯ × Singularity} {\href{https://github.com/MaximovIlya/singularity}{(repo)}}
\begin{itemize}
    \itemsep -3pt {}
     \item Полнофункциональная DeFi-платформа, где кредиторы получают доход от предоставления ликвидности, а заёмщики — доступ к кредитам под залог криптоактивов, с автоматическим управлением рисками и ставками.
\end{itemize}

\item \textbf{EasyA × Algorand London Hackathon}
\begin{itemize}
    \itemsep -3pt {}
     \item Оптимизированная платформа для выпуска NFT-коллекций, позволяющая создателям эффективно развертывать токены с кастомизируемыми параметрами при минимальных затратах и максимальной скорости.
\end{itemize}

\end{rSection}

\end{document}
